\documentclass[10pt]{article}
\usepackage{precalculus-article}
\usepackage{precalculus-key}
\newcommand{\TallMath}[1]{$\displaystyle #1\rule[-1.4em]{0pt}{3.2em}$}

\begin{document}

\pageheader{Unit 4, Lesson 4.5}{Warm-Up --- Answer Key}
\namedateperiod

\vspace{0.12in}

\begin{notesbox}{Warm-Up --- Key}
\begin{enumerate}[leftmargin=1.5em, itemsep=6pt]

\item Solve $x^2 + y^2 = 100$ for $y$. How many solutions for $y$ are there when $x = 6$?

\vspace{0.4in}
$y = $ \ans{$\pm\sqrt{100 - x^2}$} \hfill Number of solutions: \ans{2}
\vspace{0.1in}

\item The table below shows $(x,\, y)$ pairs on a curve.

\vspace{4pt}
\renewcommand{\arraystretch}{1.3}
\begin{tabular}{c|c}
$x$ & $y$ \\\hline
$1$ & $3$ \\
$2$ & $5$ \\
$4$ & $9$ \\
\end{tabular}

\vspace{6pt}
Compute $\Delta y / \Delta x$ from $x = 1$ to $x = 2$: \ans{$(5-3)/(2-1) = 2$}

Is this ratio positive or negative? \ans{Positive}

What does the sign tell you about how $y$ changes as $x$ increases? \ans{As $x$ increases, $y$ also increases (same direction).}

\vspace{0.1in}

\item The equation is $y^2 = x + 4$.
\begin{enumerate}[leftmargin=1.2em, itemsep=4pt]
    \item[(a)] For $x = 0$, what are the $y$-values? \ans{$y = \pm 2$}
    \item[(b)] For $x = 5$, what are the $y$-values? \ans{$y = \pm 3$}
\end{enumerate}

\end{enumerate}
\end{notesbox}

\begin{teachernote}
Q1: Students must write $\pm$; a common error is writing only the positive root. Emphasize that both are valid solutions.\\
Q2: The ratio $\Delta y/\Delta x = 2$ is positive, meaning $x$ and $y$ increase together. Use this to preview co-variation vocabulary in the lesson.\\
Q3: Remind students to use the square-root property. For $x = 5$: $y^2 = 9$, so $y = \pm 3$.
\end{teachernote}

\end{document}
